\documentclass[11pt,a4paper]{article}
\usepackage[T1]{fontenc}
\usepackage{lmodern,amsmath,amssymb,amsthm,booktabs,microtype}
\usepackage[margin=27mm]{geometry}
\usepackage[colorlinks=true,linkcolor=blue,citecolor=blue,urlcolor=blue]{hyperref}
\setlength{\emergencystretch}{3em}
\newtheorem{theorem}{Theorem}[section]
\newtheorem{lemma}[theorem]{Lemma}
\newtheorem{proposition}[theorem]{Proposition}
\newcommand{\Q}{\mathbb Q}
\newcommand{\Tr}{\operatorname{Tr}}
\title{Uniform complement guards on full auxiliary elliptic points}
\author{}\date{17 November 2026}
\begin{document}\maketitle
\begin{abstract}
For either even-valuation derivative detector, the completeness choice
$u\ge10$ gives $0<\eta=28w/V^2<147/242$ at every point character
of an actual guarded full auxiliary tuple. The original physical choice
$z>a+10$ already implies this weight bound for both contact sets.
Thus the finite complement equation $c^2=1-\eta$ has positive unit
radicand in every character. If its selected value is a rational square,
it supplies a native pole witness with no retained pole parameters and
nonzero terminal ordinate, using the same nine-coordinate finite source.
The required selected rational square is not proved to exist.
A constant-coefficient normalizer gives a rational complement at the
reference trace two and preserves the complete eleven-coordinate source
under its stated arithmetic premises; its actual square condition remains open.
\end{abstract}

\section{The actual point characters}
Retain the complete source and all point, ordinate, difference, quotient
and physical guards of the even-valuation construction \cite{Even}.
Its curve and root point are
\[
 K=\Q(\sqrt2),\qquad
 E:y^2=x\bigl(x^2-(3-\sqrt2)^2\bigr),\qquad A=[n]P.
\]
The root retains its square-abscissa chart, while $H,U$ use the full
guarded affine charts. This note does not restrict $H,U$ to $2E(K)$.
At an actual rational tuple, the three rank-four point blocks have
$64$ geometric characters. They are obtained by choosing the two
ordinate signs independently in the two real embeddings for each
point role. All these characters are real, since the actual ordinates
give nonzero real square roots in both embeddings. The retained
guards keep every required difference and quotient defined at every
character \cite{Full}.

The six tests are the actual source functions
\[
 \begin{gathered}
 f(Q)=1/x(Q),\quad \zeta(Q)=-x(Q)/y(Q),\qquad
 r=\zeta(U)/\zeta(H),\quad
 r'=\overline{\zeta(U-A)/\zeta(H-P)},\\
 (C_1,\ldots,C_6)=
 (f(H),f(U),\overline{f(H-P)},\overline{f(U-A)},r-r',r-X).
 \end{gathered}
\]
The last two tests may vanish. For a rational physical weight $u>0$,
common to every point character, put
\begin{equation}\label{eq:detector}
 \begin{gathered}
 D_i=\begin{cases}C_i/(1+uC_i),&1\le i\le4,\\
 C_i^2/(1+uC_i^2),&5\le i\le6,
 \end{cases}\\
 W=\prod_{i=1}^6(1+D_i^2/u),
 \quad V=\Tr(W^{-1}),\quad w=\frac{u}{(1+u)^2},\quad
 \eta=\frac{28w}{V^2}.
 \end{gathered}
\end{equation}
The scale and weight depend on the retained root abscissa and slope,
not on the ordinate signs. Thus the same rational $u$ occurs at all
characters. No independent trace or freely assigned test is used.

\begin{lemma}\label{lem:abscissa}
At either real embedding, every finite real point of $E$ has
$x> -5$. Under the retained guards all four reciprocal abscissa
tests in \eqref{eq:detector} therefore satisfy
\[
              D_i=\frac1{x(Q)+u},\qquad 0<D_i<\frac2u
                         \quad(u\ge10),
\]
with the appropriate embedding for the conjugated tests.
\end{lemma}
\begin{proof}
The two positive real curve parameters are $3-\sqrt2$ and
$3+\sqrt2$. If $x<-(3+\sqrt2)$, then
$x(x^2-(3\mp\sqrt2)^2)<0$, which cannot be a real square.
Hence $x\ge-(3+\sqrt2)>-5$. The point and difference guards
already exclude the torsion and zero-ordinate cases required to
define the tests; the bound itself also holds for those finite
points. Since $x\ne0$, the first balancing simplifies to
$1/(x+u)$. For $u\ge10$,
$x+u>u-5\ge u/2>0$, proving both assertions. The original
denominator $1+uC=(x+u)/x$ is nonzero; when $x<0$ its sign is
negative. The positivity used here is that of the cleared denominator
$x+u$ and of $D_i$.
\end{proof}

\begin{theorem}\label{thm:uniform}
At every point character of an actual guarded tuple, $u\ge10$ implies
\begin{equation}\label{eq:bound}
 1<W_\pm<(1+4/u^3)^6,\qquad
 \frac2{(1+4/u^3)^6}<V<2,\qquad
                  0<\eta<\frac{147}{242}<1.
\end{equation}
In particular all balanced denominators, detector products, traces,
and the complement $1-\eta$ are units at every character.
\end{theorem}
\begin{proof}
Lemma~\ref{lem:abscissa} treats the first four factors and their
denominators. At every real character the remaining tests are real,
so $1+uC_i^2>0$ and $0\le D_i<1/u$ for $i=5,6$.
Every detector factor is positive and smaller than $1+4/u^3$.
The first four are strictly greater than one. This proves the two
product bounds and, after taking reciprocal sums, the bounds on $V$.
Consequently
\[
 0<\eta<7\frac{u}{(1+u)^2}(1+4/u^3)^{12}
 \le \frac{70}{121}\left(\frac{251}{250}\right)^{12}
 <\frac{70}{121}\frac{21}{20}=\frac{147}{242}.
\]
Both factors in the middle expression decrease for $u\ge10$.
The final strict inequality is the exact integer comparison
\[
 21\cdot250^{12}-20\cdot251^{12}
                    =1108398992054805009292439980>0.
\]
All $64$ point characters are real as specified above, so nonzero
values at these characters prove the claimed finite-algebra units.
\end{proof}

The choice $u\ge10$ is made only in completeness. It is not imposed
as an inequality in the global sound source. All inherited guards
remain in that source; their nonvanishing here is proved, not presumed.

\section{The original physical real bound suffices}
For $k=7$ or $9$ use the retained contact sets
\[
 \mathcal S_7=\{-7,-6,-4,1,3,5,8\},\qquad
 \mathcal S_9=\{-10,-8,-5,1,2,3,4,6,7\}.
\]
With the respective fixed scales $a=M^8$ and $a=M^4$, put
\[
 F(a,z)=\prod_{c\in\mathcal S_k}((a+c)^2-z^2),\qquad
                         u=-2kF/F_a,
\]
where differentiation precedes specialization. In actual source tuples
$a>\max_c|c|$. For example $B\ge3$ gives $a\ge B^{16}>10$
in the nine-contact case, and an even larger lower bound in the
seven-contact case \cite{Even,Seven}.

\begin{proposition}\label{prop:physical}
For either contact set, $a>\max_c|c|$ and $z>a+10$ imply $u>10$.
Thus the original real choice in completeness already supplies
Theorem~\ref{thm:uniform}, uniformly in all auxiliary points.
\end{proposition}
\begin{proof}
All $a+c$ are positive and smaller than $z$. Logarithmic
differentiation gives the exact positive expression
\[
 u=\frac{2k}{\displaystyle\sum_{c\in\mathcal S_k}
                     \frac{2(a+c)}{z^2-(a+c)^2}}.
\]
Its summands satisfy
\[
 \frac{2(a+c)}{z^2-(a+c)^2}
 =\frac{2(a+c)}{z+a+c}\frac1{z-a-c}
 <\frac1{z-a-c}<\frac1{10-c}.
\]
The exact harmonic sums therefore give
\[
\begin{array}{crr}
\toprule
k&\displaystyle\sum_{c\in\mathcal S_k}\frac1{10-c}
&\text{strict lower bound for }u\\[2pt]
\midrule
7&98251/85680&1199520/98251>10\\
9&1093/840&15120/1093>10\\
\bottomrule
\end{array}
\]
The same formula proves $F,F_a,u\ne0$ on this real domain.
Also $z>a$ makes the signed height argument $a^2-z^2$ nonzero.
\end{proof}

\section{A finite complement root and the precise remaining question}
The native radicand is
\[
 P_d(v,s)=d(d+v^2)^2s^4-d(d-v^2)^2(1-d)^2.
\]
If $c^2=1-d$, the polynomial identity
\begin{equation}\label{eq:special}
                     P_d(1,c)=[2d(1-d)]^2
\end{equation}
gives a native witness with no retained pole parameters. Its ordinate
is nonzero when $d(1-d)\ne0$. This is not the degenerate witness
with zero final ordinate.

\begin{proposition}\label{prop:finite}
Replace the two free pole coordinates in the derivative source by the
single finite equation $c^2=1-\eta$. For every sufficiently large fixed
odd $B$, under ordinary rational $abc$ and the retained arithmetic
premises, this gives a sound
nine-coordinate finite source, including the input. If an actual
guarded tuple with $z>a+10$ has its selected
value $1-\eta$ in $\Q^2$, every new finite-character unit and real-root
condition holds automatically. Its native witness is
\[
                      (v,s,Y)=(1,c,2\eta(1-\eta)).
\]
No completeness assertion for the selected rational square is made.
\end{proposition}
\begin{proof}
For soundness no real bound is needed. If a discrete prime valuation
of $\eta$ is negative and odd, then
$v_p(1-\eta)=v_p(\eta)$ is also odd, contradicting its being a
rational square. Outside the retained exceptional set, which contains
seven, this excludes negative odd valuations of $b=\eta/7$.
The unchanged local extraction and height proof of \cite{Even}
therefore applies. Equivalently \eqref{eq:special} gives the native
pole equation, whose local implication permits zero terminal ordinates
\cite{Intrinsic}.

For the completeness statement in this proposition,
Theorem~\ref{thm:uniform} and Proposition~\ref{prop:physical} give
$0<\eta<147/242$ at every point character. Thus $1-\eta$ is
a positive unit everywhere; both of its roots at every real character
are nonzero and real. Also $\eta(1-\eta)$ is a unit, so the native
ordinate in the displayed witness is nonzero at every character.
The stipulated rational square supplies the rational root at the
selected character. The point rank stays $4^3=64$ and one monic
quadratic extension gives rank $128$, just as in the previous source.
The same finite-algebra encoder is shared with the point choices.
The coordinate census is
\[
                 1\text{ input}+6\text{ point coordinates}
                    +1\text{ slope}+1\text{ encoder}=9.
\]
No new free coordinate is used to impose these unit conditions.
\end{proof}

For a fixed nonzero integer input $n$, fixed defining parameter and
$A=[n]P$, the unresolved arithmetic assertion is the existence of
actual guarded $H,U$ and rational $z>a+10$ such that
\[
                         V(H,U,z)^2-28w(z)\in\Q^2.
\]
This is equivalent to the selected complement-square condition, since
$V\ne0$. The trace and weight here are the actual source functions.
The uniform positive bound and real roots do not imply that this
rational number is a rational square. The retained prime preparation
and finite local approximation statements do not supply that missing
global assertion. No nine-variable completeness, new polynomial degree,
or twenty-one-variable construction follows from this note alone.

\section{A normalizer with a rational reference complement}
Keep $u,D_i,W,V$ unchanged and define
\begin{equation}\label{eq:star}
 w_* = \frac{16u}{(4+7u)^2},\qquad b_* = \frac{4w_*}{V^2},
 \qquad \eta_* = 7b_*.
\end{equation}
Retain all old guards, including $1+u\ne0$, and add $4+7u\ne0$.
The latter is automatic on the chosen real completeness domain.

\begin{proposition}\label{prop:star}
With the intrinsic pole equation at $\eta_*$, the variant
\eqref{eq:star} retains a complete eleven-coordinate source for every
sufficiently large fixed odd $B$, under ordinary rational $abc$ and
the retained arithmetic premises. At every actual point character,
$u\ge10$ gives
\[
                0<\eta_*<\frac{294}{1369}<1.
\]
Its finite complement replacement is a sound nine-coordinate source,
with the same automatic character-unit conclusion conditional on a
selected rational square. No new degree bound is claimed.
\end{proposition}
\begin{proof}
For $p\notin\{2,7\}$, writing $\nu=v_p(u)$ gives
\[
 v_p(w_*)=\begin{cases}
 \nu,&\nu>0,\\ -\nu,&\nu<0,\\
 -2v_p(4+7u),&\nu=0.
 \end{cases}
\]
Thus positive odd contact extraction is unchanged. At negative odd
$\nu$, the unchanged balancing gives $W\equiv1$ in both slots
and $V\equiv2$, so $v_p(b_*)=-\nu>0$. At even $\nu$ the
valuation of $b_*$ is even, including primes newly dividing $4+7u$.
At selected positive odd contacts the old unit-trace preparation gives
$v_p(b_*)=\nu$. These cases retain the square-denominator proof
outside two and seven. The exact square ratio
$w_*/w=[4(1+u)/(4+7u)]^2$ alone would not prove that assertion.

At fixed places take the same sufficiently deep large-slope preparation.
Since $a$ is a rational square,
\[
 u\sim z^2/a,\qquad w_*\sim(4\sqrt a/(7z))^2,\qquad V\longrightarrow2.
\]
At two, take $\nu=v_2(u)$ negative and even with $u$ a local square.
Then $v_2(b_*)=4-\nu\ge4$ is even and $b_*$ is a local square.
At seven choose even $\nu=v_7(u)\le-4$; the unit trace gives
$v_7(b_*)=-\nu-2\ge2$. Hence $\eta_*$ has square reduced
denominator, belongs to $7\Q_2^{\times2}$, and has the required
two-adic depth. The new denominator is nonzero in these fixed
neighborhoods. No new prime outside the retained exceptional set is used.

For the real bound use Theorem~\ref{thm:uniform}'s unchanged trace estimate.
The function $16u/(4+7u)^2$ decreases for $u\ge10$, so
\[
 0<\eta_*<7w_*(1+4/u^3)^{12}
 \le\frac{280}{1369}\left(\frac{251}{250}\right)^{12}
 <\frac{294}{1369}<1.
\]
This holds on the full auxiliary domain at all point characters.
The retained intrinsic-pole theorem now gives rational witness pairs
and the same density for finite-character guard avoidance
\cite{Intrinsic}. The point roles, prime preparation, height argument,
finite ranks and retained coordinate count do not change.
Replacing those two pole coordinates by $c^2=1-\eta_*$ uses
Proposition~\ref{prop:finite} and identity~\eqref{eq:special} exactly
as before; its actual selected square is an additional requirement.
\end{proof}

The choice of constants in \eqref{eq:star} gives the exact reference identity
\[
 \left.1-\frac{28w_*}{V^2}\right|_{V=2}
 =1-7w_* = \left(\frac{4-7u}{4+7u}\right)^2.
\]
Every finite guarded actual tuple on the real domain has $V<2$ by
Theorem~\ref{thm:uniform}. Thus the identity at the reference value
does not supply an actual witness. The unresolved actual equation is
\[
                         S^2=V^2-\frac{448u}{(4+7u)^2}.
\]

The standard-library checker in \texttt{uniform\_complement\_guard\_v77},
under \nolinkurl{computation/pole_parameters/}, verifies the exact integer
margin, both harmonic sums, identity~\eqref{eq:special}, the new
normalizer identity and bound, and the previously audited source census
and guard metadata. It does not compile a new coefficient source or
search for either selected square.

\begin{thebibliography}{9}\raggedright
\bibitem{Even}\emph{Even abscissa valuations and a smaller eleven-variable degree bound},
\nolinkurl{papers/conditional_few_variables/even_valuation_detector_20261117_v75.tex}.
\bibitem{Full}\emph{Full auxiliary point charts in the derivative construction},
\nolinkurl{papers/conditional_few_variables/full_auxiliary11_20261117_v71.tex}.
\bibitem{Seven}\emph{Seven-contact detectors and a genus-twelve actual weight curve},
\nolinkurl{papers/conditional_few_variables/seven_contact_detector_20261117_v74.tex}.
\bibitem{Intrinsic}\emph{A two-coordinate pole test of coefficient height four},
\nolinkurl{papers/pole_parameters/intrinsic_pole_20261117_v56.tex}.
\end{thebibliography}
\end{document}
